\chapter{Conclusion}
\label{ch:conclusion}

This thesis investigates whether Brazil's ProUni scholarship program generates measurable returns in the formal labor market, using a novel identification strategy that combines spatial variation in program intensity with age-cohort contrasts inspired by \citeA{Duflo2001}. By comparing wages of young workers (ages 22--28, plausibly exposed to ProUni) against older workers (ages 29--35, too old to have benefited) within the same microregion and year, the analysis isolates the age-specific human capital channel from confounding local economic trends. The headline finding is a positive and statistically significant effect of ProUni intensity on the cohort wage gap in low-dose microregions, with an implied return to a ProUni-induced degree of approximately [FILL]\%. High-dose microregions exhibit attenuated returns, consistent with the saturation hypothesis. Built-in placebo tests --- using the old cohort, which should be unaffected --- confirm that the treatment effect is age-specific.

This work makes three contributions to the literature. First, it provides the first application of the Duflo-style age-cohort design to ProUni, demonstrating that the program's labor market effects are identifiable even without individual-level linked data. Second, it documents a clear saturation pattern in the returns to higher-education expansion, with policy implications for the optimal geographic allocation of scholarships. Third, it examines heterogeneity by gender and race, connecting ProUni's affirmative action provisions to differential labor market returns. The analysis employs the heterogeneity-robust estimator of \citeA{CallawayAndSantAnna2021}, buttressed by extensive robustness checks including TWFE benchmarks, doubly-robust estimation, HonestDiD sensitivity analysis, alternative discretization cutoffs, and continuous dose-response estimation.

Several avenues for future research emerge. First, linked employer-employee data (if access can be obtained) would allow individual-level tracking of ProUni recipients, sharpening the estimates and enabling analysis of occupational sorting and firm-level effects. Second, the interaction between ProUni, FIES, and Lei de Cotas deserves systematic investigation, as these programs likely interact in ways that ecological designs cannot fully capture. Third, longer-run follow-up (beyond 2019) would reveal whether the wage returns persist, grow, or decay as treated cohorts accumulate experience. Finally, extending the analysis to the informal sector --- perhaps through linked PNAD Cont\'{i}nua data --- would address the most important limitation of the RAIS-based design.
